problem:  
Let $\{(M_i^n,g_i)\}$ be a sequence of closed Riemannian $n$‑manifolds with sectional curvature bounded below by $-1$ and diameter uniformly bounded above by $1$, converging in the Gromov–Hausdorff sense to a compact smooth manifold $(X^k,h)$ of lower dimension $k<n$. By Yamaguchi’s fibration theorem, after passing to a subsequence, each $M_i$ admits a continuous surjection $\pi_i:M_i\to X$ that is almost a Riemannian submersion with almost totally geodesic fiber of small diameter and almost nilpotent structure group.

**Question:**  
Under only a lower sectional curvature bound, can these local nilpotent fibrations always be promoted—after finite covers and sufficiently small perturbations—to a genuine smooth fibration whose structure group is a compact Lie group acting isometrically on the fibers? In other words, does collapsing with a lower curvature bound and smooth limit necessarily organize itself, perhaps after resolving singularities and passing to covers, into a global smooth geometric fibration with compact structure group?

Why is it a "good" problem:  
This problem strengthens the known collapsing theory that currently requires two‑sided sectional curvature bounds for the existence of global geometric fibrations. The question probes whether analogous structure persists under *only* a lower curvature bound, pushing the boundary between Alexandrov geometry and smooth Riemannian geometry. It confronts the fundamental gap between local nilpotent (or $F$‑structure) descriptions and global group actions, asking for precise conditions transforming local collapsing models into global geometric fibrations. Resolving it would clarify the geometric nature of lower curvature bound collapses, connect Alexandrov space theory with Lie group and topological fibration theory, and potentially lead to a unified framework encompassing both smooth and singular collapsing. It is conceptually simple but demands profound insight into the interaction of curvature, topology, and group actions—hallmarks of a “good” research-level problem.